
现在要在
$\mathcal M^{\mathrm{main}}_{k+1,m}(L;\beta)$
上把仓西结构整体化。
第~7.1.2、7.1.3 小节已给出所需的分析工具；
还要使用下面的定理~7.1.44，
其证明在第~2.1 节留待此处。

\begin{theorem}{7.1.44}\label{thm:7-1-44}
参见 \textbf{[Sep84, Sep91, SeSi89]}。
若 $k+2m-2\ge0$ 且 $k\ge0$，
则第~2 节定义的
$\mathcal M^{\mathrm{main}}_{k+1,m}$
具有维数 $k+2m-2$ 的带角光滑流形结构。
\end{theorem}

\noindent\textbf{证明。}
对 $m$ 归纳。
$m=0$ 时，若 $k\ge2$，
则
$\mathcal M^{\mathrm{main}}_{k+1,0}
\cong D^{k-2}$，
即 $(k-2)$ 维圆盘；见第~2.1 节。
若 $k=1$，归纳须从
$\mathcal M^{\mathrm{main}}_{2,1}$ 开始，
它显然同胚于区间 $[0,1]$。
若 $k=0$，首个须研究的情形是
$\mathcal M^{\mathrm{main}}_{1,1}$，
它是一个点。

现讨论归纳步骤。
设 $k+2m-2\ge0$、$k\ge0$。
如下构造遗忘映射\liftIndex{I230}
\[
\operatorname{Forget}\colon
\mathcal M^{\mathrm{main}}_{k+1,m+1}
\longrightarrow
\mathcal M^{\mathrm{main}}_{k+1,m}.
\]
设
$(\Sigma,\vec z,\vec z^{\, +})$
代表
$\mathcal M^{\mathrm{main}}_{k+1,m+1}$
中的一个元素；
其中
$\vec z=(z_0,\ldots,z_k)$
是边界标记点，
$\vec z^{\, +}=(z_1^+,\ldots,z_{m+1}^+)$
是内部标记点。
记
$\vec z^{\, + -}=(z_1^+,\ldots,z_m^+)$。
若
$(\Sigma,\vec z,\vec z^{\, + -})$
稳定，则定义
\[
\operatorname{Forget}(\Sigma,\vec z,\vec z^{\, +})
=(\Sigma,\vec z,\vec z^{\, + -}).
\]
否则，(7.1.41.1)--(7.1.41.6)
中的一种情形发生；
照前述方式收缩不稳定分支，
得到
$\operatorname{Forget}(\Sigma,\vec z,\vec z^{\, +})$。
这实际上是当 $(M,L)$ 都为点时
先前的
$\operatorname{Forget}_{\rho,\rho'}$
的特例。

下一个引理描述
$\operatorname{Forget}$ 的纤维。
对任意
$(\Sigma',\vec z',\vec z^{\, +\prime})
\in\mathcal M^{\mathrm{main}}_{k+1,m}$，
如下定义只依赖于 $\Sigma'$ 的
$\widetilde\Sigma'$：
把 $\Sigma'$ 的每个边界奇点
变形成光滑点，
但不变形内部奇点。
因此，$\widetilde\Sigma'$
是带有若干球面气泡的一个圆盘。
这里只考虑它的 $C^\infty$ 结构，
不考虑复结构；
在此意义下
$\widetilde\Sigma'$ 得到良定义。

\begin{lemma}{7.1.45}\label{lem:7-1-45}
遗忘映射
$\operatorname{Forget}\colon
\mathcal M^{\mathrm{main}}_{k+1,m+1}
\to\mathcal M^{\mathrm{main}}_{k+1,m}$
在
$(\Sigma',\vec z',\vec z^{\, +\prime})$
处的纤维微分同胚于
$\widetilde\Sigma'$。
\end{lemma}

在 \textbf{[FuOn99II]} 第~23 节，
针对带标记点的闭黎曼曲面有类似论证。

\begin{proof}
设
$\operatorname{Forget}(\Sigma,\vec z,\vec z^{\, +})
=(\Sigma',\vec z',\vec z^{\, +\prime})$。
沿用遗忘映射定义中的记号，
分情形讨论。

\noindent\textbf{情形 1。}
若
$(\Sigma,\vec z,\vec z^{\, + -})$
稳定，
则它等于
$(\Sigma',\vec z',\vec z^{\, +\prime})$。
此时
$(\Sigma,\vec z,\vec z^{\, +})$
与
$\Sigma'\setminus
(\partial\Sigma'\cup\vec z^{\, +\prime}
\cup\{\text{奇点}\})$
中的点一一对应。

\noindent\textbf{情形 2。}
若发生 (7.1.41.1)，
容易看出此情形不可能出现。

\noindent\textbf{情形 3。}
若发生 (7.1.41.2)，
则这样的
$(\Sigma,\vec z,\vec z^{\, +})$
与 $\Sigma'$ 的内部奇点一一对应；
见图~\ref{fig:7-1-14}。

\noindent\textbf{情形 4。}
若发生 (7.1.41.3)，
则这样的
$(\Sigma,\vec z,\vec z^{\, +})$
与 $\vec z^{\, +\prime}$ 中的点，
即已有的内部标记点，一一对应；
见图~\ref{fig:7-1-15}。

\noindent\textbf{情形 5。}
若发生 (7.1.41.4)，
则这样的
$(\Sigma,\vec z,\vec z^{\, +})$
与
$\partial\Sigma'\setminus
(\vec z'\cup\{\text{奇点}\})$
中的点一一对应；
见图~\ref{fig:7-1-16}。

\noindent\textbf{情形 6。}
若发生 (7.1.41.5)，
则对应于 $\Sigma'$ 的一个边界奇点。
但对每个这样的边界奇点，
存在一族由
$\mathcal M^{\mathrm{main}}_{2,1}$
参数化的
$(\Sigma,\vec z,\vec z^{\, +})$；
这一参数空间是一个区间。
见图~\ref{fig:7-1-17}。

\noindent\textbf{情形 7。}
若发生 (7.1.41.6)，
则这样的
$(\Sigma,\vec z,\vec z^{\, +})$
与 $\vec z'$ 中的点，
即已有的边界标记点，一一对应；
见图~\ref{fig:7-1-18}。

\par\medskip\noindent
\begin{minipage}[t]{0.32\textwidth}\centering
\setlength{\unitlength}{0.70mm}
\begin{picture}(50,61)
\put(12,45){\circle{12}}
\put(12,31){\circle{11}}
\put(12,16){\circle{18}}
\put(12,39){\circle*{1}}
\put(12,25){\circle*{1}}
\put(12,31){$\scriptstyle\times$}
\put(20,34){$\scriptstyle z_{m+1}^+$}
\put(23,27){\vector(1,0){5}}
\put(39,31){\circle{12}}
\put(39,16){\circle{18}}
\put(39,25){\circle*{1}}
\end{picture}
\par{\renewcommand{\thefigure}{7.1.14}\refstepcounter{figure}
\normalfont\bfseries 图~\thefigure\label{fig:7-1-14}}
\end{minipage}\hfill
\begin{minipage}[t]{0.32\textwidth}\centering
\setlength{\unitlength}{0.70mm}
\begin{picture}(50,61)
\put(12,45){\circle{12}}
\put(12,31){\circle{11}}
\put(12,16){\circle{18}}
\put(12,39){\circle*{1}}
\put(12,25){\circle*{1}}
\put(9,47){$\scriptstyle\times$}
\put(13,43){$\scriptstyle\times$}
\put(19,49){$\scriptstyle z_i^+$}
\put(18,39){$\scriptstyle z_{m+1}^+$}
\put(23,27){\vector(1,0){5}}
\put(39,31){\circle{11}}
\put(39,16){\circle{18}}
\put(39,25){\circle*{1}}
\put(39,31){$\scriptstyle\times$}
\put(42,36){$\scriptstyle z_i^+$}
\end{picture}
\par{\renewcommand{\thefigure}{7.1.15}\refstepcounter{figure}
\normalfont\bfseries 图~\thefigure\label{fig:7-1-15}}
\end{minipage}\hfill
\begin{minipage}[t]{0.32\textwidth}\centering
\setlength{\unitlength}{0.70mm}
\begin{picture}(50,61)
\put(12,38){\circle{13}}
\put(12,18){\circle{19}}
\put(12,28){\circle*{1}}
\put(12,38){$\scriptstyle\times$}
\put(17,42){$\scriptstyle z_{m+1}^+$}
\put(23,27){\vector(1,0){5}}
\put(39,18){\circle{19}}
\end{picture}
\par{\renewcommand{\thefigure}{7.1.16}\refstepcounter{figure}
\normalfont\bfseries 图~\thefigure\label{fig:7-1-16}}
\end{minipage}\par\medskip

\noindent
\begin{minipage}[t]{0.48\textwidth}\centering
\setlength{\unitlength}{0.68mm}
\begin{picture}(72,66)
\put(14,56){\circle{8}}
\put(14,46){\circle{10}}
\put(14,33){\circle{12}}
\put(14,17){\circle{18}}
\put(14,52){\circle*{1}}
\put(14,40){\circle*{1}}
\put(14,26){\circle*{1}}
\put(14,33){$\scriptstyle\times$}
\put(19,35){$\scriptstyle z_{m+1}^+$}
\put(30,32){\vector(1,0){8}}
\put(54,50){\circle{8}}
\put(43,42){\circle{8}}
\put(54,38){\circle{11}}
\put(54,22){\circle{18}}
\put(54,44){\circle*{1}}
\put(54,30){\circle*{1}}
\put(47,45){\line(1,-1){3}}
\put(35,9){$\scriptstyle\mathcal M^{\mathrm{main}}_{2,1}\cong[0,1]$}
\end{picture}
\par{\renewcommand{\thefigure}{7.1.17}\refstepcounter{figure}
\normalfont\bfseries 图~\thefigure\label{fig:7-1-17}}
\end{minipage}\hfill
\begin{minipage}[t]{0.48\textwidth}\centering
\setlength{\unitlength}{0.68mm}
\begin{picture}(72,66)
\put(14,46){\circle{13}}
\put(14,23){\circle{23}}
\put(14,35){\circle*{1}}
\put(14,46){$\scriptstyle\times$}
\put(19,50){$\scriptstyle z_{m+1}^+$}
\put(7,54){$\scriptstyle z_2$}
\put(26,30){$\scriptstyle z_1$}
\put(26,20){$\scriptstyle z_0$}
\put(3,20){$\scriptstyle z_3$}
\put(7,10){$\scriptstyle z_4$}
\put(30,23){\vector(1,0){8}}
\put(53,23){\circle{23}}
\put(47,35){$\scriptstyle z_2$}
\put(63,30){$\scriptstyle z_1$}
\put(63,20){$\scriptstyle z_0$}
\put(40,20){$\scriptstyle z_3$}
\put(45,10){$\scriptstyle z_4$}
\end{picture}
\par{\renewcommand{\thefigure}{7.1.18}\refstepcounter{figure}
\normalfont\bfseries 图~\thefigure\label{fig:7-1-18}}
\end{minipage}\par\medskip

综上，纤维由 $\Sigma'$
将每个边界奇点替换为区间 $[0,1]$
而得到；这正是上文定义的
$\widetilde\Sigma'$。
引理~\ref{lem:7-1-45} 得证。
\end{proof}

\noindent\textbf{定理~\ref{thm:7-1-44} 的证明（续）。}
由引理~\ref{lem:7-1-45}，
现在可以完成归纳步骤。
假设
$\mathcal M^{\mathrm{main}}_{k+1,m}$
是带角光滑流形，
要证明
$\mathcal M^{\mathrm{main}}_{k+1,m+1}$
亦然。
设
$(\Sigma,\vec z,\vec z^{\, +})
\in\mathcal M^{\mathrm{main}}_{k+1,m+1}$，
且
\[
\operatorname{Forget}
(\Sigma,\vec z,\vec z^{\, +})
=(\Sigma',\vec z',\vec z^{\, +\prime}).
\]
按该引理，
$(\Sigma,\vec z,\vec z^{\, +})$
对应于
$\widetilde\Sigma'$ 中一点 $x$。
若 $x$ 不是奇点，
由归纳假设容易在其附近构造坐标图；
即使 $x$ 是边界点也一样。
故设 $x$ 是奇点。
按定义，
$\widetilde\Sigma'$
已没有边界奇点，
所以 $x$ 必为内部点。
于是可以使用无边界黎曼曲面模空间的同样论证，
在
$(\Sigma,\vec z,\vec z^{\, +})$
附近寻找坐标图。

更具体地说，
在
$(\Sigma',\vec z',\vec z^{\, +\prime})$
于
$\mathcal M^{\mathrm{main}}_{k+1,m}$
的某个邻域中，
存在复余维一的除子 $S$；
只有 $S$ 上的元素才使 $x$ 保持为奇点。
由归纳假设，
底空间的这一邻域可与
\[
\mathbb R_{\ge0}^{d'}\times
\mathbb R^{d-d'}\times\mathbb C
\]
中零点的邻域相识别，
并且
\[
S=\mathbb R_{\ge0}^{d'}\times
\mathbb R^{d-d'}\times\{0\},
\qquad
d=k+1+2m-5
=\dim\mathcal M^{\mathrm{main}}_{k+1,m}-2.
\]
在这一识别下取坐标
$(u_1,\ldots,u_d,q)$，
其中每个 $u_i$ 属于
$\mathbb R$ 或 $\mathbb R_{\ge0}$，
$q\in\mathbb C$。
此时
$(\Sigma,\vec z,\vec z^{\, +})$
在
$\mathcal M^{\mathrm{main}}_{k+1,m+1}$
中的邻域可与
\[
\mathbb R_{\ge0}^{d'}\times
\mathbb R^{d-d'}\times\mathbb C^2
\]
中零点的邻域相识别，
其中遗忘映射为
\[
(u_1,\ldots,u_d,q_1,q_2)
\longmapsto
(u_1,\ldots,u_d,q_1q_2).
\]
事实上，
在
$S=\{(u_1,\ldots,u_d,q)\mid q=0\}$
上一点的纤维与上述邻域的交集由
$q_1q_2=0$ 给出；
它在 $q_1=q_2=0$ 处有一个
与 $x$ 对应的节点。
定理~\ref{thm:7-1-44} 得证。\hfill$\square$

\begin{remark}{7.1.46}\label{rem:7-1-46}
由定理~\ref{thm:7-1-44}
的证明及引理~\ref{lem:7-1-45}，
可以类似于
F. F. Knudsen \textbf{[Knu83]} 和
S. Keel \textbf{[Kee92]}
研究亏格零带标记点稳定曲线模空间那样，
描述
$\mathcal M^{\mathrm{main}}_{k+1,m}$
的同调群。
例如，当 $k+1\ge3$ 时，
该空间在 $m=0,1,2$ 时可缩，
但在 $m=3$ 时有非平凡同调。
其中最简单的是
\[
H_2(\mathcal M^{\mathrm{main}}_{3,3};\mathbb Q)
\cong
H^4(\mathcal M^{\mathrm{main}}_{3,3},
\partial\mathcal M^{\mathrm{main}}_{3,3};\mathbb Q)
\]
中的非平凡同调类。
它由一个集合的闭包的基本类给出。
该集合的元素为
\[
(\Sigma,\vec z,(z_1^+,z_2^+,z_3^+)),
\]
其中 $\Sigma$ 是一个圆盘，
在 $0\in D^2$ 处附有一个球面气泡，
而 $z_1^+,z_2^+,z_3^+$
都位于该气泡上。
容易看出，
这一集合是六维流形
$\mathcal M^{\mathrm{main}}_{3,3}$
中的二维子流形。

可以仿照通常的格罗莫夫--威腾不变量，
将第~3.8 节算子
$\mathfrak q$ 的序列构造
与
$\mathcal M^{\mathrm{main}}_{k+1,m}$
的同调类耦合而作推广。
此事留待今后研究。
\end{remark}

现在可以完成命题~\ref{prop:7-1-1}、
\ref{prop:7-1-2} 的证明。
为此须把迄今构造的仓西邻域拼接为
整体仓西结构。
利用
$\mathcal M^{\mathrm{main}}_{k+1}(\beta)$
的紧性与豪斯多夫性，
只需拼接有限多个仓西邻域。
此拼接的细节与
\textbf{[FuOn99II]} 第~4.2 节相同；
下面回顾其中的论证。
对每个
$\sigma=((\Sigma,\vec z),w)
\in\mathcal M^{\mathrm{main}}_{k+1}(\beta)$，
已经选取了有限维子空间
\[
E'_\sigma\subset
\Gamma\bigl(
\Sigma,\Lambda^{0,1}_\Sigma\otimes w^*TM
\bigr),
\]
使方程
$\bar\partial w'\equiv0\pmod{E'_\sigma}$
定义 $\sigma$ 的仓西邻域
$U'_\sigma$。
前一小节把 $E'_\sigma$
记作
$E_{((\mathfrak t,\ell),w,\vec p)}$。
附录 A1 用 $V$ 表示仓西邻域；
这里改用 $U$，
因为前两个小节中的 $V$ 已另有所指。
适当缩小后，
可假设 $U'_\sigma$ 是 $\sigma$ 的闭邻域。
由紧性，选取有限多个点
$\sigma_1,\ldots,\sigma_N$，使
\begin{equation}
\bigcup_{i=1}^{N}
\operatorname{Int}U'_{\sigma_i}
\supset
\mathcal M^{\mathrm{main}}_{k+1}(\beta).
\tag{7.1.47}\label{eq:7-1-47}
\end{equation}
记
$\sigma_i=((\Sigma_i,\vec z(i)),w_i)$。
利用覆盖~\eqref{eq:7-1-47}，
将修改 $E'_\sigma$
来构造 $\sigma$ 的仓西邻域。

先设 $(\Sigma,\vec z)$ 稳定。
收集所有满足
$\sigma\in U'_{\sigma_i}$
的 $i$，并令
\[
E_\sigma=
\bigoplus_{i:\,\sigma\in U'_{\sigma_i}}
E'_{\sigma_i}.
\]
与本小节前半类似，
将各
$E'_{\sigma_i}$
视为
$\Gamma(\Sigma,\Lambda^{0,1}_\Sigma\otimes w^*TM)$
的子空间。
仓西邻域 $U_\sigma$
是满足下列条件的同构类所成之集：
对
$((\Sigma',\vec z'),w')$，
其域
$(\Sigma',\vec z')$
在定义~2.1.16 给出的
$\mathcal M^{\mathrm{main}}_{k+1}$
拓扑下接近
$(\Sigma,\vec z)$，
且
$\bar\partial w'\equiv0\pmod{E_\sigma}$。

由于
$U'_{\sigma_i}$ 闭，
可选择 $U_\sigma$，
使
\[
\sigma_1\in U_{\sigma_2}
\quad\Longrightarrow\quad
E_{\sigma_1}\subseteq E_{\sigma_2}.
\]
这里借助由平行移动定义的映射
\[
\Gamma(\Sigma_1,\Lambda^{0,1}\otimes w_1^*TM)
\longrightarrow
\Gamma(\Sigma_2,\Lambda^{0,1}\otimes w_2^*TM)
\]
来比较这两个有限维空间，
其中
$\sigma_i=((\Sigma_i,\vec z_i),w_i)$。
于是还可假设：
若 $\sigma_1\in U_{\sigma_2}$，
则
$U_{\sigma_1}$ 中 $\sigma_1$
的一个邻域包含在
$U_{\sigma_2}$ 内。
将该邻域记为
$U_{\sigma_1,\sigma_2}$；
它满足附录 A1 的条件 (A1.12)。

若 $(\Sigma,\vec z)$ 不稳定，
则使用定理~\ref{thm:7-1-44}
如下处理。
以下论证是
\textbf{[FuOn99II]} 附录中论证的复述；
也可以采取类似该书第~15 节的另一种方法。
增添若干内部标记点
$\vec z^{\, +}$，
使
$(\Sigma,\vec z,\vec z^{\, +})$
稳定。
还可以不失一般性地假设：
\begin{enumerate}[label=\textup{(7.1.48.\arabic*)},leftmargin=6em]
\item\label{eq:7-1-48-1}
$\vec z^{\, +}$ 的每一点 $z_i^+$
都位于 $w$ 非平凡的分支上；
\item\label{eq:7-1-48-2}
$w$ 在 $\vec z^{\, +}$ 的每一点处都是浸入。
\end{enumerate}
不失一般性的理由与
\textbf{[FuOn99II]} 附录相同。
对每个
$z_i^+\in\vec z^{\, +}$，
取 $M$ 中余维为 $2$ 的子流形
$Q_i^{2n-2}$，
使它与 $w(\Sigma)$
在 $w(z_i^+)$ 处横截相交。
利用定理~\ref{thm:7-1-44}，
可仿照前述 $m=0$ 的情形，
在
$\mathcal M_{k+1,m}(\beta)$
中为
$((\Sigma,\vec z,\vec z^{\, +}),w)$
构造仓西邻域；
这里 $m$ 是
$\vec z^{\, +}$ 的点数。
内部标记点处的评价映射为
\[
\operatorname{ev}_{\mathrm{int}}\colon
\mathcal M_{k+1,m}(\beta)\longrightarrow M^m.
\]
可假设此映射弱浸没。
于是下列空间带有仓西结构：
\begin{equation}
\mathcal M_{k+1,m}(\beta)
{}_{\operatorname{ev}_{\mathrm{int}}}
\!\times_{M^m}
\prod_{i=1}^{m}Q_i.
\tag{7.1.49}\label{eq:7-1-49}
\end{equation}
由 (7.1.48.2) 容易看出，
式~\eqref{eq:7-1-49} 的空间中一点
\[
\bigl(
((\Sigma,\vec z,\vec z^{\, +}),w),
(w(z_1^+),\ldots,w(z_m^+))
\bigr)
\]
的仓西邻域，
可与
$\mathcal M^{\mathrm{main}}_{k+1}(\beta)$
中
$((\Sigma,\vec z),w)$
的仓西邻域相识别。
故即使
$(\Sigma,\vec z)$
不稳定，
也已经为点
$((\Sigma,\vec z),w)$
构造了仓西邻域。
余下论证与域稳定时相同。
命题~\ref{prop:7-1-1}、
\ref{prop:7-1-2} 证毕。\hfill$\square$
